\documentclass{article}
\usepackage{amsmath, amssymb, amsthm}

\setlength{\parindent}{0pt}

\newtheorem{claim}{Claim}

\begin{document}

There are about $250{,}000{,}000$ people in the United States who might use a phone. To keep the problem 
simple, we ignore the fact that people are on the phone more often at certain times of day and on certain 
days of the year. Assume that each person is on the phone during each minute mutually independently with 
probability $p = 0.01$.

\begin{claim}
    The expected number of people on the phone at a given moment is $2{,}500{,}000$.
\end{claim}

\begin{proof}
    For each person $i$, define the indicator random variable
    \[
    X_i =
    \begin{cases}
    1 & \text{if person } i \text{ is on the phone} \\
    0 & \text{otherwise}
    \end{cases}
    \]
    Since a person is on the phone with probability $0.01$,
    \[
    \mathrm{E}[X_i] = 1(0.01) + 0(0.99) = 0.01
    \]

    \medskip

    Let $X$ denote the total number of people on the phone at a given moment. Then
    \[
    X = X_1 + X_2 + \dots + X_{250{,}000{,}000}
    \]
    By linearity of expectation,
    \[
    \begin{aligned}
    \mathrm{E}[X] &= \sum_{i=1}^{250{,}000{,}000}\mathrm{E}[X_i] \\ 
    &= 250{,}000{,}000(0.01) \\
    &= 2{,}500{,}000
    \end{aligned}
    \]

    \medskip

    Hence the expected number of people on the phone at a given moment is $2{,}500{,}000$.
\end{proof}

\begin{claim}
    Suppose that we construct a phone network whose capacity is a mere one percent above the expectation.
    Then the probability that the network is overloaded in a given minute is at most 
    $2^{-165.64} \approx 1.37 \times 10^{-50}$. % not the tightest (checked answer key) but good enough
\end{claim}

\begin{proof}
    Let $X$ denote the number of people on the phone at a given moment, so $X \sim \mathrm{Binomial}(n, p)$
    with $n = 250{,}000{,}000$ and $p = 0.01$. By Claim $1$, $\mathrm{E}[X] = 2{,}500{,}000$.

    \medskip
    
    Let $a = 1.01 \cdot \mathrm{E}[X] = 2{,}525{,}000$ denote the network's capacity. The network is
    overloaded precisely when $X > a$. Since $X$ is integer-valued, this is equivalent to
    \[
    X \geq a + 1 = 2{,}525{,}001
    \]

    \medskip

    We may write the tail probability as
    \[
    \Pr\{X > a\} = \sum_{k = a+1}^{n} \mathrm{PDF}_X(k)
    \]
    Since $a+1$ already exceeds $\mathrm{E}[X]$, the terms $\mathrm{PDF}_X(k)$ are decreasing for
    $k \geq a+1$. Hence every term in the sum is at most $\mathrm{PDF}_X(a+1)$, and the sum has
    $n - a$ terms, so
    \[
    \Pr\{X > a\} \leq (n-a)\cdot \mathrm{PDF}_X(a+1)
    \]

    \medskip

    Write $a + 1 = \alpha n$, where
    \[
    \alpha = \frac{a+1}{n} = \frac{2{,}525{,}001}{250{,}000{,}000}
    \]
    By the Stirling-based approximation of the PDF,
    \[
    \mathrm{PDF}_X(\alpha n) \leq \frac{2^{nH(\alpha)}}{\sqrt{2\pi\alpha(1-\alpha)n}}\cdot p^{\alpha n}(1-p)^{(1-\alpha)n}
    \]
    where $H(\alpha) = \alpha\log_2\frac{1}{\alpha} + (1-\alpha)\log_2\frac{1}{1-\alpha}$.

    \medskip

    Putting the two bounds together,
    \[
    \Pr\{X > a\} \leq (n-a)\cdot \frac{2^{nH(\alpha)}}{\sqrt{2\pi\alpha(1-\alpha)n}}\cdot p^{\alpha n}(1-p)^{(1-\alpha)n}
    \]
    Since the right-hand side involves numbers far too large to evaluate directly, we take
    $\log_2$ of both sides (this is valid since $\log_2$ is increasing):
    \[
    \begin{aligned}
    \log_2\Pr\{X > a\} \leq \log_2(n-a) + nH(\alpha) - \frac{1}{2}\log_2\left(2\pi\alpha(1-\alpha)n\right) \\ 
    +\; \alpha n\log_2(p) + (1-\alpha)n\log_2(1-p)
    \end{aligned}
    \]

    \medskip

    Substituting $n = 250{,}000{,}000$, $\alpha = \frac{2{,}525{,}001}{250{,}000{,}000}$, and $p = 0.01$,
    \[
    \begin{aligned}
    \log_2\Pr\{X > a\} &\leq \log_2\left(250{,}000{,}000-2{,}525{,}001\right) \\
    &+250{,}000{,}000H\left(\frac{2{,}525{,}001}{250{,}000{,}000}\right) \\
    &-\frac{1}{2}\log_2\left(2\pi\left(\frac{2{,}525{,}001}{250{,}000{,}000}\right)
    \left(1-\frac{2{,}525{,}001}{250{,}000{,}000}\right)250{,}000{,}000\right) \\
    &+\frac{2{,}525{,}001}{250{,}000{,}000}250{,}000{,}000\log_2(0.01) \\
    &+ \left(1-\frac{2{,}525{,}001}{250{,}000{,}000}\right)250{,}000{,}000\log_2(0.99) \\
    &\approx -165.64
    \end{aligned}
    \]

    \medskip

    Hence the probability that the network is overloaded in a given minute is at most
    \[
    \Pr\{X > a\} \leq 2^{-165.64} \approx 1.37 \times 10^{-50}
    \]
\end{proof}

\begin{claim}
    The expected number of minutes until the system is overloaded for the first time is approximately 
    $7.29 \times 10^{49}$.
\end{claim}

\begin{proof}
    By Claim $2$, the probability of the network overloading in a given moment is at most $2^{-165.64}$. 
    Thus the expected number of minutes until the system is overloaded for the first time is 
    \[
    \frac{1}{2^{-165.64}} \approx 7.29 \times 10^{49}
    \]

    \medskip

    Hence the expected number of minutes until the system is overloaded for the first time is approximately
    $7.29 \times 10^{49}$.
\end{proof}

\end{document}
